The proof is the same as that of 1.2.10.
\begin{proposition}{6.7.2}\label{XXI.6.7.2}
Let $\Delta$ be a system of simple roots. Put
\[
E_\Delta(\mathcal R)=\{u\in\Aut(\mathcal R)\mid u(\Delta)=\Delta\}.
\]
Then $\Aut(\mathcal R)$ is the semidirect product of $W(\mathcal R)$ by $E_\Delta(\mathcal R)$.
\end{proposition}
This follows at once from the fact that $W(\mathcal R)$ acts simply transitively on the systems of simple roots, and that if $\Delta$ is a system of simple roots of $R$, $u(\Delta)$ is a system of simple roots for every automorphism $u$ of $\mathcal R$.

We shall see later a simpler description of $E_\Delta(\mathcal R)$ in the case of reduced irreducible root data.{\renewcommand{\thefootnote}{(27)}\nde{When $\mathcal R$ is simply connected or adjoint, see 7.4.5.}}
\begin{definition}{6.7.3}\label{XXI.6.7.3}
Denote by $\Aut^s(\mathcal R)${\renewcommand{\thefootnote}{(28)}\nde{The exponent $s$ is intended to suggest ``semisimple''.}} the set of $u\in\Aut(\mathcal R)$ such that the following diagram commutes:
\[
\xymatrix{\mathcal R\ar[rr]^u\ar[dr]&&\mathcal R\ar[dl]\\
&\operatorname{rad}(\mathcal R)&}.
\]
Put $E_\Delta^s(\mathcal R)=E_\Delta(\mathcal R)\cap\Aut^s(\mathcal R)$.
\end{definition}
\begin{remark}{6.7.4}\label{XXI.6.7.4}
If $u\in\Aut(\mathcal R)$, one therefore has $u\in\Aut^s(\mathcal R)$ if and only if $(u-\mathrm{id})(M)\subset\mathcal V(R)$. In particular, $W(\mathcal R)\subset\Aut^s(\mathcal R)$. It follows at once that:
\end{remark}
\begin{proposition}{6.7.5}\label{XXI.6.7.5}
The group $\Aut^s(\mathcal R)$ is the semidirect product of $W(\mathcal R)$ by $E_\Delta^s(\mathcal R)$, for every system of simple roots $\Delta$.
\end{proposition}
With every automorphism of $\mathcal R$ there is associated by functoriality an automorphism of $\operatorname{ad}(\mathcal R)$. One therefore has a canonical morphism
\[
\Aut(\mathcal R)\to\Aut(\operatorname{ad}(\mathcal R)).
\]
\begin{lemma}{6.7.6}\label{XXI.6.7.6}
The morphism $\Aut^s(\mathcal R)\to\Aut(\operatorname{ad}(\mathcal R))$ is injective.
\end{lemma}
Indeed, let $u$ be an automorphism of $M$ such that $(u-\mathrm{id})(M)\subset\mathcal V(R)$ and ${}^tu(\alpha^*)=\alpha^*$ for $\alpha^*\in R^*$. For every $x\in M$, one has
\[
(\alpha^*,u(x)-x)=({}^tu(\alpha^*)-\alpha^*,x)=0,
\]
hence $u(x)-x=0$ by 1.2.5.
\begin{lemma}{6.7.7}\label{XXI.6.7.7}
The group $\Aut^s(\mathcal R)$ is finite.
\end{lemma}
Indeed, it suffices to prove that $\Aut(\mathcal R)$ is finite if $\mathcal R$ is adjoint. Since $M$ is then generated by $R$, every automorphism of $\mathcal R$ is determined by the permutation of $R$ which it defines.
\begin{remark}{6.7.8}\label{XXI.6.7.8}
One sees at once that $\Aut(\mathcal R)$ (resp. $E_\Delta(\mathcal R)$) is finite if and only if $\operatorname{rgred}(\mathcal R)-\operatorname{rgss}(\mathcal R)\leq1$.
\end{remark}

\sgasection{6.8}{$p$-morphisms of reduced root data}
In this section, $p$ is an integer $>0$, fixed once and for all.
\begin{definition}{6.8.1}\label{XXI.6.8.1}
Let $\mathcal R=(M,M^*,R,R^*)$ and $\mathcal R=(M',M'^*,R',R'^*)$ be two reduced root data. One says that a morphism of groups
\[
f:M'\to M
\]
is a \emph{$p$-morphism} from $\mathcal R'$ to $\mathcal R$ if the following conditions hold: there exist a bijection
\[
u:R\isomto R'
\]
and a map $q:R\to\{p^n,n\in\NN\}$ such that:
\begin{enumeratei}
\item $f(u(\alpha))=q(\alpha)\alpha$ for every $\alpha\in R$.
\item ${}^tf(\alpha^*)=q(\alpha)u(\alpha)^*$ for every $\alpha\in R$.{\renewcommand{\thefootnote}{(29)}\nde{Note that these two conditions imply $q(\alpha)(u(\alpha^*),u(\beta))=q(\beta)(\alpha^*,\beta)$ for all $\alpha,\beta\in R$.}}
\end{enumeratei}
% typo?: missing prime on the second datum and u(alpha-star) in note (29) kept as printed.
\end{definition}
\begin{corollary}{6.8.2}\label{XXI.6.8.2}
A $1$-morphism is precisely a morphism.
\end{corollary}
\begin{corollary}{6.8.3}\label{XXI.6.8.3}
The transpose of a $p$-morphism is a $p$-morphism.
\end{corollary}
\begin{lemma}{6.8.4}\label{XXI.6.8.4}
If $w\in W(\mathcal R)$, $\alpha\in R$, one has $q(w(\alpha))=q(\alpha)$. The map $s_\alpha\mto s_{u(\alpha)}$ extends to an isomorphism $\overline u:W(\mathcal R)\to W(\mathcal R')$ such that
\[
u(w(\alpha))=\overline u(w)(u(\alpha)).
\]
\end{lemma}
It suffices to prove that for $\alpha,\beta\in R$, one has $u(s_\alpha(\beta))=s_{u(\alpha)}u(\beta)$ and $q(s_\alpha(\beta))=q(\beta)$. Now one has successively:
\[
\begin{aligned}
f(s_{u(\alpha)}u(\beta))&=f(u(\beta))-(u(\alpha)^*,u(\beta))f(u(\alpha))\\
&=q(\beta)\beta-q(\beta)q(\alpha)^{-1}(\alpha^*,\beta)q(\alpha)\alpha\\
&=q(\beta)(\beta-(\alpha^*,\beta)\alpha)=q(\beta)s_\alpha(\beta).
\end{aligned}
\]
If $\gamma=u^{-1}(s_{u(\alpha)}u(\beta))$, one therefore has $q(\gamma)\gamma=f(u(\gamma))=q(\beta)s_\alpha(\beta)$. The two roots $\gamma$ and $s_\alpha(\beta)$ are therefore proportional (over $\QQ$), hence equal or opposite, but $q(\gamma)$ and $q(\beta)$ are positive. One therefore has $q(\gamma)=q(\beta)$ and $\gamma=s_\alpha(\beta)$.
\begin{definition}{6.8.5}\label{XXI.6.8.5}
The $q(\alpha)$ are called the \emph{root exponents} of $f$.
\end{definition}
\begin{example}{6.8.6}\label{XXI.6.8.6}
Let $\mathcal R$ be a reduced root datum and $q=p^n$ ($n\in\NN$). Then multiplication by $q:M\to M$, $x\mto qx$, is a $p$-morphism all of whose root exponents are $q$ (and $u=\mathrm{id}$); it is denoted
\[
q:\mathcal R\to\mathcal R.
\]
\end{example}
\begin{proposition}{6.8.7}\label{XXI.6.8.7}
With the notations of 6.8.1, $u$ gives an isomorphism from the set of systems of simple roots (resp. positive roots) of $R$ onto the corresponding set for $R'$.
\end{proposition}
This follows at once from 3.1.5 (resp. 3.2.1).

\sgasection{7}{Structure}\label{XXI.sec.7}
\sgasection{7.1}{Decomposition of a root datum}
\begin{proposition}{7.1.1}\label{XXI.7.1.1}
Let $\mathcal R$ be a root datum, $\Delta$ a system of simple roots.
\begin{enumeratei}
\item Let $R'$ and $R''$ be two closed symmetric sets of roots forming a partition of $R$. If one puts $\Delta'=\Delta\cap R'$, $\Delta''=\Delta\cap R''$, then $R'=R_{\Delta'}$, $R''=R_{\Delta''}$, and every root of $\Delta'$ is orthogonal to every root of $\Delta''$.
\item Let $\Delta'$ and $\Delta''$ be two orthogonal subsets of $\Delta$ forming a partition of $\Delta$. Then $R'=R_{\Delta'}$ and $R''=R_{\Delta''}$ form a partition of $R$.
\end{enumeratei}
\end{proposition}
Let us first prove (i).
\begin{lemma}{7.1.2}\label{XXI.7.1.2}
Under the conditions of (i), if $\alpha$, $\beta$ and $\alpha+\beta$ are roots, all three belong to $R'$ or all three to $R''$.
\end{lemma}
Suppose, for example, $\alpha+\beta\in R'$. Then one cannot have $\alpha,\beta\in R''$, since $R''$ is closed; suppose therefore $\alpha\in R'$. Then $-\alpha\in R'$, and $\beta=(\beta+\alpha)-\alpha\in R'$.

Let us now show that $R'=R_{\Delta'}$ by induction on the order of a positive root $\alpha\in R'\cap\mathcal P(\Delta)$. If $\operatorname{ord}_\Delta(\alpha)=1$, then $\alpha\in R'\cap\Delta=\Delta'$. If $\operatorname{ord}_\Delta(\alpha)>1$, there exists $\beta\in\Delta$ such that $\alpha-\beta\in R$. By the lemma, one has $\beta\in\Delta'$, $\alpha-\beta\in R'$, hence $\alpha-\beta\in R_{\Delta'}$ by induction, and finally $\alpha=(\alpha-\beta)+\beta\in R_{\Delta'}$.

Finally, let us show that $\Delta'$ and $\Delta''$ are orthogonal. If $\alpha\in\Delta'$ and $\beta\in\Delta''$, then $(\beta^*,\alpha)\leq0$. If $(\beta^*,\alpha)\ne0$, $\beta+\alpha$ is a root, contrary to the lemma.

Let us prove (ii). If $\Delta'$ or $\Delta''$ is empty, this is immediate. Otherwise, if $R_{\Delta'}$ and $R_{\Delta''}$ do not form a partition of $R$, there exists a root $\alpha$ of the form
\[
\alpha=\sum m'_i\alpha'_i+\sum m''_j\alpha''_j,\qquad m'_i\in\ZZ_+,\quad m''_j\in\ZZ_+,
\]
where $\alpha'_i$ (resp. $\alpha''_j$) denote elements of $\Delta'$ (resp. $\Delta''$). Applying 3.1.2, one deduces a relation of the form (interchanging $\Delta'$ and $\Delta''$ if necessary):
\[
\delta=\gamma+\beta,\qquad\gamma\in R_{\Delta'},\quad\beta\in\Delta'',\quad\delta\in R.
\]
But since $(\beta^*,\gamma)=0$, $\gamma-\beta$ is also a root by 2.2.5, which is impossible.
\begin{proposition}{7.1.3}\label{XXI.7.1.3}
Let $\mathcal R$ be a root datum. The following conditions are equivalent:
\begin{enumeratei}
\item There exists no nontrivial partition of $R$ into two closed symmetric subsets.
\item For one (resp. every) system of simple roots $\Delta$ of $R$, there exists no partition of $\Delta$ into two nonempty orthogonal subsets.
\item The natural representation of $W(\mathcal R)$ on $\mathcal V(R)$ is irreducible.
\item For every pair $(\alpha,\beta)$ of roots, there exists a sequence of roots $\alpha_0,\alpha_1,\alpha_2,\ldots,\alpha_n$, with $\alpha=\alpha_0$, $\alpha_n=\beta$, the roots $\alpha_i$ and $\alpha_{i+1}$ ($i=0,\ldots,n-1$) being nonorthogonal.
\end{enumeratei}
\end{proposition}
One has (i) $\Leftrightarrow$ (ii) by 7.1.1. Clearly (iv) $\Rightarrow$ (ii). Conversely, if (ii) holds for $\Delta$, condition (iv) holds if $\alpha,\beta\in\Delta$. Now, for every root there exists a simple root not orthogonal to it (3.1.1, for example). On the other hand, (iii) $\Rightarrow$ (i). Indeed, under the conditions of 7.1.1, $\mathcal V(R')$ is stable under $W(\mathcal R)$. It remains to prove (i) $\Rightarrow$ (iii).

Let therefore $H$ be a vector subspace of $\mathcal V(R)$ stable under $W(\mathcal R)$. For every $\alpha\in R$, the equation $s_\alpha(H)=H$ gives at once $\alpha\in H$ or $\alpha^*\in H^\perp$ (the orthogonal of $H$ in $\mathcal V(R^*)$, which is in duality with $\mathcal V(R)$). Putting $R'=\{\alpha\in R\mid\alpha\in H\}$ and $R''=\{\alpha\in R\mid\alpha^*\in H^\perp\}$, one has obtained a partition of $R$ into two closed symmetric subsets.
\begin{definition}{7.1.4}\label{XXI.7.1.4}
A root datum (resp. root system) satisfying the equivalent conditions of 7.1.3 and of semisimple rank $\ne0$ is called \emph{irreducible}.
\end{definition}
\begin{corollary}{7.1.5}\label{XXI.7.1.5}
For every root datum $\mathcal R$, there exists a unique partition (up to order) of $R$ into closed symmetric irreducible subsets.
\end{corollary}
\begin{corollary}{7.1.6}\label{XXI.7.1.6}
Every adjoint (resp. simply connected) root datum is a product of irreducible adjoint (resp. simply connected) root data.
\end{corollary}
It suffices to see this in the adjoint case. The assertion then follows from the fact that under the conditions of 7.1.1 one has
\[
\Gamma_0(R)=\Gamma_0(R')\times\Gamma_0(R'').
\]
\begin{corollary}{7.1.7}\label{XXI.7.1.7}
For every root datum (resp. reduced root datum) $\mathcal R$, there exists an isogeny $\mathcal R\to\mathcal R'$, where $\mathcal R'$ is a product of a trivial root datum and irreducible (resp. reduced) simply connected root data.
\end{corollary}

\sgasection{7.2}{Properties of irreducible root data}
\begin{definition}{7.2.1}\label{XXI.7.2.1}
Let $\mathcal R$ be an irreducible root datum. For every $\alpha\in R$, put
\[
\operatorname{long}(\alpha)=\ell(\alpha)/\ell(\alpha_0);
\]
where $\alpha_0\in R$ is such that $\ell(\alpha_0)$ is minimal; one says that $\operatorname{long}(\alpha)$ is the \emph{length} of $\alpha$.
\end{definition}
\begin{lemma}{7.2.2}\label{XXI.7.2.2}
Let $\mathcal R$ be an irreducible root datum. The Weyl group acts transitively on the set of roots of the same length.
\end{lemma}
Indeed, let $\alpha,\beta\in R$. Since the representation of $W$ on $\mathcal V(R)$ is irreducible, $\alpha$ cannot be orthogonal to all the $w(\beta)$, $w\in W$. There therefore exists $w\in W$ with $w(\beta)$ nonorthogonal to $\alpha$. Now $\ell(w(\beta))=\ell(\beta)$, and one concludes by 2.3.2.
\begin{lemma}{7.2.3}\label{XXI.7.2.3}
If $\mathcal R$ is irreducible and reduced, $\operatorname{long}(R)$ is $\{1\}$, $\{1,2\}$ or $\{1,3\}$.
\end{lemma}
By the remark used above, for every $\alpha,\beta\in R$ there always exists $w\in W$ such that $w(\beta)$ is not orthogonal to $\beta$. One therefore has $\ell(\alpha)/\ell(\beta)=1$, $2$, $3$, $1/2$ or $1/3$ (by 2.3.1). One therefore has $\operatorname{long}(\alpha)=1$, $2$ or $3$, but if $\operatorname{long}(\alpha)=2$, $\operatorname{long}(\beta)=3$, then $\ell(\alpha)/\ell(\beta)=2/3$, which is impossible.
% typo?: beta rather than alpha as the orthogonality target kept as printed.
\begin{remark}{7.2.4}\label{XXI.7.2.4}
Reasoning similarly, one proves the following result: if $R$ is irreducible and nonreduced with $\operatorname{rgss}(\mathcal R)>1$, one has $\operatorname{long}(R)=\{1,2,4\}$. Putting $\operatorname{long}^{-1}(i)=R_i$, one has $\operatorname{ind}(R)=R_1\cup R_2$, $R_4=2R_1$, and two nonproportional roots of $R_1$ are orthogonal. Conversely, if $R$ is an irreducible reduced system such that $\operatorname{long}(R)=\{1,2\}$, put $\operatorname{long}^{-1}(i)=R_i$ and suppose that two nonproportional roots of $R_1$ are orthogonal; then $R\cup2R_1$ is irreducible, nonreduced, and $\operatorname{ind}(R\cup2R_1)=R$.
\end{remark}
\begin{lemma}{7.2.5}\label{XXI.7.2.5}
If $\mathcal R$ is an irreducible root datum, $\mathcal R^*$ is also one, and the product
\[\operatorname{long}(\alpha)\operatorname{long}(\alpha^*)\]
is constant as $\alpha$ runs through $R$.
\end{lemma}
This follows at once from 7.1.3 (iv) and 2.2.6.
\begin{definition}{7.2.6}\label{XXI.7.2.6}
Let $\mathcal R$ be any root datum. One calls the \emph{length} of $\alpha\in R$, and denotes by $\operatorname{long}(\alpha)$, the length of $\alpha$ in its irreducible component.
\end{definition}
\begin{lemma}{7.2.7}\label{XXI.7.2.7}
There exists a unique group homomorphism $u:\Gamma_0(R)\to\Gamma_0(R^*)$ such that $u(\alpha)=\operatorname{long}(\alpha)\alpha^*$ for $\alpha\in R$.
\end{lemma}
By 3.5.5, it suffices to check that if $\alpha,\beta,\alpha+\beta\in R$, one has
\[
\operatorname{long}(\alpha)\alpha^*+\operatorname{long}(\beta)\beta^*
=\operatorname{long}(\alpha+\beta)(\alpha+\beta)^*.
\]
But $\alpha$, $\beta$ and $\alpha+\beta$ are in the same irreducible component of $R$ by 7.1.2, and one is reduced to 1.2.2.
\begin{remark}{7.2.8}\label{XXI.7.2.8}
Let $u$ be as in 7.2.7. For $\alpha,\beta\in R$, one has $(u(\alpha),\beta)=(u(\beta),\alpha)$. Indeed, this amounts to seeing that
\[
\operatorname{long}(\alpha)(\alpha^*,\beta)=\operatorname{long}(\beta)(\beta^*,\alpha),
\]
which clearly holds if $\alpha$ and $\beta$ are orthogonal. If $\alpha$ and $\beta$ are not orthogonal, they are in the same irreducible component of $R$, and one is reduced to 1.2.1, formula (9).
\end{remark}
\begin{remark}{7.2.9}\label{XXI.7.2.9}
The symmetric bilinear form $(u(x),y)$ is positive nondegenerate on $\Gamma_0(R)$.
\end{remark}
Indeed, let $R_i$ be the irreducible components of $R$. One has
\[
\Gamma_0(R)=\prod_i\Gamma_0(R_i),
\]
and the bilinear form $(u(x),y)$ is the product of the forms
\[
2^{-1}\ell(\alpha_i)^{-1}(p(x),y)
\]
on the $\Gamma_0(R_i)$, where $\ell(\alpha_i)$ is the minimum of $\ell(\alpha)$ for $\alpha\in R_i$. Now these various symmetric bilinear forms are positive nondegenerate (1.2.6).

\sgasection{7.3}{Cartan matrix}
Let $\mathcal R$ be a root datum. If $\Delta$ is a system of simple roots, one calls the \emph{Cartan matrix} of $\mathcal R$ relative to $\Delta$ the square matrix indexed by $\Delta$ defined by
\[
a_{\alpha,\beta}=(\alpha^*,\beta)\quad\text{for }\alpha,\beta\in\Delta.
\]
Note first that if $\Delta'$ is another system of simple roots and $w$ an element of $W(\mathcal R)$ such that $w(\Delta)=\Delta'$, one has
\[
(w(\alpha)^*,w(\beta))=(\alpha^*,\beta),
\]
hence the Cartan matrix of $\mathcal R$ relative to $\Delta'$ is obtained from that relative to $\Delta$ by the isomorphism $\Delta\to\Delta'$ of index sets defined by $w$. It follows that, up to an isomorphism of the index set, the Cartan matrix depends only on $\mathcal R$.
\begin{proposition}{7.3.1}\label{XXI.7.3.1}
The Cartan matrix has the following properties:
\begin{enumeratei}
\item $a_{\alpha,\alpha}=2$,\qquad$a_{\alpha,\beta}\leq0$ for $\alpha\ne\beta$.
\item $a_{\alpha,\beta}=0$ implies $a_{\beta,\alpha}=0$.
\item There exist strictly positive integers $m_\alpha$ ($=\operatorname{long}(\alpha)$) such that the matrix $(m_\alpha a_{\alpha,\beta})$ is symmetric, positive and nondegenerate.
\item The diagonal minors of the matrix $(a_{\alpha,\beta})_{\alpha,\beta\in\Delta}$, i.e. the determinants
\[
\det(a_{\alpha,\beta})_{\alpha,\beta\in\Delta'}\quad\text{for }\Delta'\subset\Delta,
\]
are strictly positive.
\item One has $s_\alpha(\beta)=\beta-a_{\alpha,\beta}\alpha$ and $s_\alpha(\beta^*)=\beta^*-a_{\beta,\alpha}\alpha^*$.
\end{enumeratei}
\end{proposition}
Indeed, (v) is a definition, (i) follows from 3.2.11, (ii) from 2.2.2, (iii) from 7.2.9, and (iv) follows at once from (iii) by the relation
\[
\det(m_\alpha a_{\alpha,\beta})_{\alpha,\beta\in\Delta'}=
\left(\prod_{\alpha\in\Delta'}m_\alpha\right)\det(a_{\alpha,\beta})_{\alpha,\beta\in\Delta'}.
\]
\begin{proposition}{7.3.2}\label{XXI.7.3.2}
Let $\mathcal R$ and $\mathcal R'$ be two simply connected (resp. adjoint) reduced root data, $\Delta$ (resp. $\Delta'$) a system of simple roots of $\mathcal R$ (resp. $\mathcal R'$), and $u:\Delta\to\Delta'$ an isomorphism such that if $(a_{\alpha,\beta})$ and $(a'_{\alpha',\beta'})$ denote the Cartan matrices of $\mathcal R$ and $\mathcal R'$ relative to $\Delta$ and $\Delta'$, one has:
\[
a'_{u(\alpha),u(\beta)}=a_{\alpha,\beta}.
\]
Then there exists a unique isomorphism from $\mathcal R$ onto $\mathcal R'$ inducing $u$ on $\Delta$.
\end{proposition}
Clearly, it suffices to give the proof in the adjoint case. Then $M=\Gamma_0(R)$ and $M'=\Gamma_0(R')$ are the free abelian groups generated by $\Delta$ and $\Delta'$. There therefore exists a unique group isomorphism from $M$ onto $M'$ inducing $u$ on $\Delta$. Denote it also by $u$. Let us show that $u(R)\subset R'$. Every root $\alpha$ of $\mathcal R$ is written $s_{\alpha_1}\cdots s_{\alpha_n}(\alpha_{n+1})$, with $\alpha_i\in\Delta$. Clearly,
\[
u(\alpha)=s_{u(\alpha_1)}\cdots s_{u(\alpha_n)}(u(\alpha_{n+1})),
\]
by the hypothesis on $u$ and relations (v) of 7.3.1.
